\chapter{Cableado del poste}\label{cap:cableado}

\begin{figure}[H]\centering
\resizebox{\linewidth}{!}{% Cableado de la PCB_V2 según el esquema (redes de KiCad en letra pequeña)
\begin{tikzpicture}[
  font=\sffamily\small,
  bornera/.style={draw=grafito, fill=lima!35, rounded corners=1pt, align=center, minimum width=27mm, minimum height=10mm},
  pieza/.style={draw=grafito, fill=marfil, rounded corners=2pt, align=center, minimum height=10mm, inner sep=4pt},
  fusible/.style={draw=grafito, fill=white, align=center, minimum width=13mm, minimum height=6mm, font=\sffamily\scriptsize},
  aviso/.style={draw=rojoaviso, fill=rojoaviso!8, align=center, minimum width=13mm, minimum height=8mm, font=\sffamily\scriptsize},
  v12/.style={-{Stealth[length=2mm]}, thick, rojoaviso!80!black},
  dat/.style={-{Stealth[length=2mm]}, thick, azulnota},
  red/.style={font=\ttfamily\scriptsize, inner sep=1.5pt},
]
  % --- Entrada 1: regulador e inyección ---
  \node[bornera] (in1) at (0,0) {S\_LM\_INY1\\+12 V · GND};
  \node[fusible] (f5) at (3.6,0) {F5};
  \node[pieza] (lm) at (7.2,0) {LM2596\\12 V → \dato{reguladorSalida}};
  \node[pieza] (esp) at (12.2,0) {ESP32-S3\\5V / GND};
  \node[fusible] (f1) at (3.6,-1.6) {F1};
  \node[fusible] (f2) at (3.6,-2.8) {F2};
  \node[bornera] (iny1) at (7.2,-1.6) {INY1 · inyección 1};
  \node[bornera] (iny2) at (7.2,-2.8) {INY2 · inyección 2};
  \draw[v12] (in1) -- node[above, red, pos=0.45]{+12V\_INY} (f5);
  \draw[v12] (f5) -- node[above, red]{+12V\_LM} (lm);
  \draw[v12] (lm) -- node[above, red]{+5V\_ESP} (esp);
  \fill[rojoaviso!80!black] (2.25,0) circle (0.6mm);
  \draw[v12] (2.25,0) |- (f1.west);
  \draw[v12] (2.25,0) |- (f2.west);
  \draw[v12] (f1) -- (iny1);
  \draw[v12] (f2) -- (iny2);

  % --- Entrada 2: caras ---
  \node[bornera] (in2) at (0,-5.4) {S\_A\_B1\\+12 V · GND};
  \node[fusible] (f6) at (3.6,-4.8) {F6};
  \node[fusible] (f7) at (3.6,-6.0) {F7};
  \node[bornera] (a1) at (7.2,-4.8) {A1 → cara A\\GND · D1 · +12V};
  \node[bornera] (b1) at (7.2,-6.0) {B1 → cara B\\GND · D2 · +12V};
  \draw[v12] (in2.east) -- ++(4mm,0) |- (f6.west);
  \draw[v12] (in2.east) -- ++(4mm,0) |- (f7.west);
  \draw[v12] (f6) -- node[above, red]{+12V\_A} (a1);
  \draw[v12] (f7) -- node[below, red]{+12V\_B} (b1);

  % --- Datos ---
  \node[pieza, align=center] (b2) at (12.2,-5.4) {Bornera B2\\GPIO\_A · GPIO\_B · GPIO\_C\\R1 → D1 · R2 → D2};
  \draw[dat] (esp.south) -- node[right, font=\sffamily\scriptsize, align=left]{señales\\de datos} (b2.north);
  \draw[dat] (b2.west) ++(0,1.5mm) -- ++(-5mm,0) |- (a1.east);
  \draw[dat] (b2.west) ++(0,-1.5mm) -- ++(-5mm,0) |- (b1.east);
  \node[font=\sffamily\scriptsize, align=center] at (12.2,-6.85) {Poste que funciona: cara A en GPIO\dato{pinesPosteActual.caraA}, cara B en GPIO\dato{pinesPosteActual.caraB}\\ESP32 soldada en la \dato{pcbVersion}: GPIO\dato{pinesEsquema.GPIO_A} y GPIO\dato{pinesEsquema.GPIO_B}};
\end{tikzpicture}
}
\caption{Conexiones de la \dato{pcbVersion} según el esquema.}\label{fig:cableado}
\end{figure}

\advertencia{La fuente se conecta a la red de 120~V. Haz todas las conexiones con la red desconectada y usa terminales adecuados. El poste se conecta con un cable dúplex de dos hilos y una clavija sin polo a tierra, así que el borne de tierra de la fuente queda libre. Por eso la fuente debe ir dentro de la caja cerrada, sin tocar partes metálicas que el público pueda alcanzar. Si la caja o el poste son metálicos, usa un cable de tres hilos y conecta la tierra. La entrada de corriente alterna debe quedar fuera del alcance del público.}

\section{Fuente y placa}

\begin{enumerate}
\item Fuente: L y N al cable dúplex de 2×18~AWG con su clavija (con la red desconectada). El borne de tierra queda libre (ver la advertencia de arriba).
\item Salida de 12~V de la fuente a S\_LM\_INY1 (borne 2 +12~V, borne 1 GND).
\item Salida de 12~V de la fuente a S\_A\_B1 (borne 2 +12~V, borne 1 GND).
\item Pon un fusible externo en el cable de la fuente que va a S\_LM\_INY1. Protege el cable entero; en la placa, F5, F1 y F2 protegen cada ramal. En una placa fabricada antes de la corrección de F5 (capítulo~\ref{cap:pcb}), este fusible externo es la única protección del regulador.
\end{enumerate}

\advertencia{La fuente \dato{fuente}. No la expongas al sol ni a la lluvia: debe ir en una caja cuya temperatura real no pase de 40~°C, o cambiarse por una fuente para exteriores. La protección IP68 de los píxeles no cubre sus empalmes, la placa ni la fuente.}

\section{Caras de luz}

Cada cara es una cadena de \dato{ledsPorCara} píxeles. Los píxeles WS2811 tienen dirección: la flecha impresa en el cable indica hacia dónde viajan los datos. El extremo de \textbf{entrada} (DIN) es el que se conecta a la placa.

\begin{enumerate}
\item Cara A a la bornera A1: borne 1 GND, borne 2 datos, borne 3 +12~V.
\item Cara B a la bornera B1, en el mismo orden.
\item Inyección: lleva INY1 y INY2 (GND y +12~V) al extremo lejano de cada cara para que el brillo sea parejo en toda la altura.
\end{enumerate}

\nota{Los colores de los cables cambian entre lotes de píxeles. Verifica con el multímetro cuál es +12~V antes de atornillar.}

\nota{La placa no lleva adaptador de nivel lógico (74HCT125): la ESP32 envía la señal a 3,3~V y funcionó en las pruebas con estos píxeles. Si con cables de datos largos aparecen parpadeos o colores errados en el primer píxel, acorta el cable de datos o añade un adaptador de nivel.}


\section{Caja y montaje}

Placa y fuente van dentro de una caja estanca, con prensaestopas para los cables. La caja debe permitir disipar el calor de la fuente y dejar accesible el USB de la ESP32 para mantenimiento.

\foto{caja-montada}{Caja abierta del poste que funciona: arriba, la placa con la ESP32-S3; abajo, la fuente de 12~V con su portafusible externo («FUSE»). Los cables rojo y negro llevan 12~V de la fuente a las borneras de la placa, y los de las caras salen por el agujero de arriba.}{0.85}
\foto{caja-cerrada}{Caja cerrada, con los cables de las dos caras saliendo por un costado y sus conectores de 3 polos hacia los píxeles WS2811.}{0.55}

\advertencia{En las fotos, los cables salen por un agujero sin prensaestopas y la caja deja de ser estanca. Antes de instalar la caja en la calle, pon un prensaestopas por cable (se compraron PG-7, PG-9 y PG-13.5; capítulo~\ref{cap:materiales}) del tamaño del cable que pasa, y sella cualquier otro agujero.}

\section{Fusibles y calibre del cable}

Los fusibles son de vidrio de 5×20~mm, los que se compraron (capítulo~\ref{cap:materiales}). Con la corriente medida (\dato{corrienteMedida}):
\begin{itemize}
\item Caras (F6 y F7) e inyección (F1 y F2): \dato{fusibleCara}. Cada cara toma \dato{corrientePorCara}, poco más de la mitad del valor del fusible.
\item Regulador (F5): \dato{fusibleRegulador}. El LM2596 con la ESP32 toma menos de 0,3~A de 12~V. Si el fusible se funde al conectar la fuente, por el pico de carga de los condensadores, cámbialo por uno de 1~A de acción lenta.
\item A la salida de la fuente, antes de dividir los ramales: \dato{fusibleFuente}. Lleva las dos caras (\dato{corrientePoste}) más el regulador, unos 3,5~A en total.
\end{itemize}
Un fusible de vidrio no debe trabajar por encima del 75~\% de su valor, y aguanta menos cuando hace calor. El de 5~A de la fuente es el más ajustado. El poste funcionó casi 4~horas seguidas con la bandera de Colombia y otros efectos sin problema; si se funde con el poste en blanco al 100~\% en un día de calor, baja el brillo máximo en LUMA o cambia ese fusible por uno de 6,3~A. Los portafusibles de la placa son abiertos: solo sirven dentro de la caja estanca.

